\documentclass{article}
\usepackage{graphicx} % Required for inserting images
\usepackage{float}
\usepackage{physics}
\usepackage{hyperref}
\usepackage{subcaption}
\usepackage{amsmath}
\usepackage{amssymb}
\usepackage[left=25mm,right=25mm]{geometry}


\begin{document}


%T�st� alkaa selkkarin kansilehti. Vaihda tilalle tarvittavat tiedot.
\begin{titlepage}
\pagestyle{empty}
\begin{center}

\vspace*{3cm}
\noindent\LARGE{\textbf{
Tietokannat Ryhmätyö
}}

%*******************************************************
\vspace{0.5cm}
\large{
\begin{tabular}{l l}
Hugo Palin & Hugo.Palin@aalto.fi\\
Elias Myllymäki & Elias.Myllymaki@aalto.fi\\
Antti Alitalo & Antti.Alitalo@aalto.fi\\
\\
\\
\end{tabular}
}
\end{center}
\end{titlepage}



\begin{figure}
    \centering
    \includegraphics[width=\linewidth]{UML-kaavio.png}
    \caption{Tietokannan UML-kaavio}
    \label{fig:UML-kaavio}
\end{figure}

% Tässä oli siis taarpeilla "Count" eli kuinka monta, mutta se poistettiin koska sitä ei itseasiassa tarvita.
\begin{verbatim}
#Alleviivauksien sijaan, Primary Key:t ovat merkitty "(PK)" merkinnöllä


# Projektit kuvailee siis rakennusfirman eri projekteja
Projektit(
    Projekti_id (PK),
    Valmis,
    Alkupvm,
    Loppupvm,
    Paikka
)

# Osaprojektit ovat yksittäisen projektin osia
# Projekti_id viittaa "Projektit"-taulukkoon
Osaprojektit(
    Osaprojekti_id (PK),
    Projekti_id,  # viittaa Projektit(Projekti_id)
    Valmis,
    Alkupvm,
    Loppupvm
)

# Kuvaa osaprojektien välisiä riippuvuuksia
OsaprojektiRiippuvuus(
    Edeltava_id (PK),  # viittaa Osaprojektit
    Seuraava_id (PK)   # viittaa Osaprojektit
)

# Työntekijät
Tyontekijat(
    Tyontekija_id (PK)
)

# Työntekijöiden pätevyydet tämä siis on vissiin erillinen rooliin verrattuna
Patevyydet(
    Patevyys_id (PK)
)

# Täytyy olla ns. "monesta moneen"-suhde (ehkä???) työntekijän ja pätevyyden välillä
TyontekijaPatevyys(
    Tyontekija_id (PK),  # viittaa Tyontekijat
    Patevyys_id (PK)     # viittaa Patevyydet
)

# Työkoneiden mallit ja yleistä tietoa niistä
Konemalli(
    Malli (PK),
    Valmistaja,
    Massa,
    Polttoaineen_kulutus,
    Kuvaus
)

# Yksittäiset työkoneet (EI siis malli vaan yksittäiset "koneyksilöt")
# Malli viittaa "Konemalli"-taulukkoon
Koneet(
    Kone_id (PK),
    Malli  # viittaa Konemalli(Malli)
)

# Yleinen tarve (tämä on siis yläluokka kaikenlaisille tarpeille)
Tarpeet(
    Tarve_id (PK)
)

# Liittää tarpeet osaprojekteihin
OsaprojektiTarve(
    Osaprojekti_id (PK),  # viittaa Osaprojektit
    Tarve_id (PK)         # viittaa Tarpeet
)

# Työkoneen tarve tietyllä aikavälillä
Konetarpeet(
    Tarve_id (PK),  # viittaa Tarpeet
    Malli,          # viittaa Konemalli
    Alkupvm,
    Loppupvm
)

# Työntekijän pätevyyteen perustuva tarve
Tyontekijatarve(
    Tarve_id (PK),   # viittaa Tarpeet
    Patevyys_id      # viittaa Patevyydet
)

# Varaus siis kuvaa resurssin käyttöä osaprojektissa jollakin aikavälillä
Varaus(
    Varaus_id (PK),
    Osaprojekti_id,  # viittaa Osaprojektit
    Alkupvm,
    Loppupvm
)

# Yksittäisen koneen varaus
Konevaraus(
    Varaus_id (PK),  # viittaa Varaus
    Kone_id          # viittaa Koneet
)

# Työntekijän varaus osaprojektiin
TyontekijaVaraus(
    Varaus_id (PK),   # viittaa Varaus
    Tyontekija_id     # viittaa Tyontekijat
)
\end{verbatim}


\newpage

% Tähän yappaamista ja selityksiä

% Anomaliat

% poistoanomaliat (engl. deletion anomalies): monikoiden poistolla voi olla sivuvaikutuksia. Jos esimerkiksi tuotteiden ja niiden valmistajien tiedot on yhdistetty samaan relaatioon, voi tuotteen poistaminen poistaa myös informaation valmistajan nimestä ja puhelinnumerosta.

Tietokannassa ei ole toisteisuutta, koska tietokanta on suunniteltu siten ettei ole toistettua informaatiota.

Osaprojekti ei voi alkaa ennen kuin projekti alkaa eikä osaprojekti voi loppua projektin loppumispäivämäärän jälkeen. Tämä on siis päivitysanomalia meidän tietokannassa.

Poistoanomalioita ei ole.
\\

% Päivitysanomalio on se, että projektin loppupäivämäärä ei voi olla ennen osaprojektin loppupäivämäärää. 
% Päivämäärät ovat päivitys anomaliat

% Jos tarvit kolme tyyppiä, niin kolme riviä lisää ??? Tarpeissa anomaly...

% Kolme riviä tarpeita, ja sitten sä etit tuolta varausoliosta että kuinka monta ihmistä joka pystyy täyttämään ne tarpeen on samaan aikaan, mutta ainoo tapa miten määrittää se on käyttämällä sellasta count funktioo muutenkin...
% Sulla on se count juttu ja voit tehä theta joinin ja varmistaa että ihmisten määrä sama kun tarpeitten määrä.

% Nyt on siis kolme riviä tietoo jokka on redundancy... (ehkä???)

% Runtimellä ei väliä tässä kurssissa???????? (sus, eli count juttu toimis sillon, vaikka joudutaan laskee joka kerta.).


Jokaisen olion primary key funktionaalisesti määrää olion muut ominaisuudet (triviaalisesti). Muita funktionaalisia riippuvuuksia ei ole
\\

Tietokanta on Boyce-Codd--normaalimuodossa

% Esimerkkejä tähän kohtaan???



\newpage

1. Jos halutaan lisätä projekteja/osaprojekteja, täytyy vain lisätä uusia projekti-olioita tai osaprojekti-olioita\quad :D
\\

2. Jos halutaan lisätä työntekijöitä johonkin osaprojektiin, niin valitaan minkälaista pätevyyttä halutaan ja sitten lisätään osaprojektiin $N$ ihmistä, kellä kyseinen pätevyys on.
\\

3. Jos halutaan lisätä jo olemassa olevia pätevyyksiä työntekijälle, tulee "TyontekijaPatevyys"-attribuuttiin lisätä rivi jossa on kyseisen työntekijän ID ja kyseisen pätevyyden ID. Jos pätevyys on kokonaan uusi, tulee kyseinen pätevyys lisätä "Patevyydet" olioon ja tämän jälkeen yhdistää työntekijä ja pätevyys yhteen edellämainitulla tavalla.
\\


% 4. Työkoneet lisäys
4. Työkonemalleja voi lisätä lisäämällä "Konemalli"-olioita joilla on tietyt parametrit. Itse työkoneinstansseja voi lisätä "Koneet"-oliolla jossa voi laittaa sarjanumeron sekä konemallin. Jokaisella työkoneella täytyy olla vastaava konemalli, mutta tietysti samaa mallia voidaan jakaa monen eri koneinstanssin kesken.
\\

5. Selvittääkseen mitä projekteja on esimerkiksi tällä hetkellä käynnissä  täytyy filtteröidä "Projektit"-oliossa olevat rivit, joiden Alkupvm on pienempi tai yhtäsuuri kuin nykyinen päivä ja jotka eivät ole valmiita. 
\\

6. Osaprojektien järjestyksen saa selville yhdistämällä "OsaprojektiRiippuvuus"-attribuutin avulla. Attribuutissa on kaikki ID:t projekteista jotka riippuvat toisistaan. Jos esimerkiksi halutaan tietää mitkä kaikki osaprojektit täytyy olla valmiita ennen osaprojekti B:tä, täytyy OsaprojektiRiippuvuudessa filteröidä kaikki rivit joissa "Seuraava\_id"==B.ID ja katsoa mitkä projektit ovat sarakkeella "Edeltava\_id . Projektilla EI saa olla samanaikaisia osaprojekteja.
\\

7. Jos haluaa tietää millä pätevyyksillä varustettuja työntekijöitä tarvitaan mihinkin osaprojektiin, täytyy tehdä uusi rivi "OsaprojektiTarve"-attribuuttiin. 
\\

8. Jos halutaan varata tietty määrä työntekijöitä johonkin osaprojektiin tietyillä pätevyyksillä, täytyy yhdistää työntekijät ja varaukset "Työntekijänvaraus" attribuutin kautta, liittää tähän tauluun kaikki pätevyydet ja filteröidä niistä toivotut pätevyydet. Tämä taulu liitetään työntekijöihin. Tällä taululla on nyt vain työntekijöitä joilla on halutut pätevyydet. Nyt taulu liitetään "Varaus"-olioon ja filteröitään kaikki työntekijät jotka eivät ole varattuja tuolloin. Lisäämällä alkioita "Varaus"-olioon ja  "Varaustarpeet"-attribuuttiin saadaan yhdistettyjä tiettyjä varauksia tiettyihin osaprojekteihin. 
\\



% Counttia ei edes kuuluisi olla tarpeet jutussa, koska joku projekti voi tarvita viis kirvesmiestä ja joku toinen kaksikymmentä...


\end{document}

